% BHU PYQ -- Computer Science IT Tools & Applications (2025-26 NEP)
\documentclass[11pt,a4paper]{article}
\usepackage[a4paper,top=2.5cm,bottom=2.5cm,left=2.8cm,right=2.8cm,headheight=14pt]{geometry}
\usepackage{amsmath,amssymb}
\usepackage[T1]{fontenc}
\usepackage[utf8]{inputenc}
\usepackage{lmodern,microtype,fancyhdr,enumitem,hyperref,lastpage,parskip}

\hypersetup{colorlinks=true,urlcolor=black,linkcolor=black,pdftitle=CSCMV32: IT Tools & Applications -- BHU NEP 2025-26}

\pagestyle{fancy}\fancyhf{}
\renewcommand{\headrulewidth}{0.4pt}\renewcommand{\footrulewidth}{0.4pt}
\fancyhead[L]{\small\textit{Banaras Hindu University}}
\fancyhead[R]{\small\textit{CSCMV32: IT Tools \& Applications (2025-26)}}
\fancyfoot[C]{\small Page \thepage\ of \pageref{LastPage}}

\newlist{parts}{enumerate}{2}
\setlist[parts,1]{label=(\alph*),leftmargin=*,itemsep=4pt,topsep=3pt}
\setlist[parts,2]{label=(\roman*),leftmargin=*,itemsep=2pt,topsep=2pt}
\newcommand{\pts}[1]{\hfill\textit{[#1]}}

\begin{document}
\begin{center}
  {\large\bfseries BANARAS HINDU UNIVERSITY}\\[2pt]
  {\normalsize B.Sc. (CS)-III Mid Semester Examination 2025-26 (NEP)}\\[6pt]
  \rule{0.8\linewidth}{0.5pt}\\[6pt]
  {\large\bfseries COMPUTER SCIENCE}\\[4pt]
  {\large\bfseries Paper: CSCMV32 --- IT Tools \& Applications}\\[4pt]
  {\normalsize Time: 3 Hours \hfill Full Marks: 40}
\end{center}
\rule{\linewidth}{0.4pt}
\smallskip
\noindent\textit{Instructions: Answer all questions as specified. The figures in the right-hand margin indicate marks.}
\bigskip

\noindent\textbf{Question 1.} Create an Excel file and name it as your exam roll number.\pts{10 Marks}
\begin{parts}
  \item In sheet1, enter 15 in A1, 25 in B1 and 30 in C1. Find the value of \texttt{=(A1+B1)/C1+2}.
  \item In sheet2, enter "Students ID" column with values \texttt{2025CSMV001} to \texttt{2025CSMV010}, add another column as "Roll Number" and enter last 3 digits from student ID using a formula. Also, write the formula for extracting roll number from student ID.
  \item Add column "Marks Obtained" and enter 10 random numbers between 1--100 in column C as the marks in English. Write the formula to enter random numbers between 1--100.
  \item Find the 2nd lowest marks obtained by the student and then use INDEX-MATCH to lookup the student ID corresponding to that number.
  \item In column "Viva Date", enter random dates between 1 December 2025 and today's date. Write the formula for the same.
  \item Using VLOOKUP, find the Viva Date of the student with Student ID = \texttt{2025CSCMV006}. Also, write the formula.
  \item Create column "Final marks" and use \texttt{LET} function along with a conditional function to enter the value equal to "Marks Obtained" $+ 3$ if "Marks Obtained" is $< 50$ and equal to "Marks Obtained" otherwise. Also, write the formula.
\end{parts}

\bigskip
\noindent\textbf{Question 2.}\pts{10 Marks}
\begin{parts}
  \item Write something about (i) Yourself (ii) Varanasi (iii) Your Hometown.
  \item Write the content in 2 different paragraphs, having at least 3 lines (not 3 sentences) each. Make 1st paragraph Centrally Aligned \& 2nd paragraph Justified. Wherever you mention your name/Varanasi/Your Hometown's Name, change the font style to Bradley Hand ITC, make it bold and underline it. Add space after paragraph \& change line spacing to 1.5 pts. Also, highlight all the content of 2nd paragraph in "Turquoise" colour with text colour in "Blue, Accent 1, Darker 25\%". Insert an image (source of image should be online) in the document according to the option chosen and format the Picture Style to RRR (Reflected Rounded Rectangle).
  \item Save the file with a filename \texttt{<Your Name> Work}. Find your name/Varanasi/Your Hometown's Name in the document, replace it with India and save these changes in a new file with filename \texttt{<Your Name> New}, in Portable Document Format (PDF).
\end{parts}

\bigskip
\noindent\textbf{Question 3.} Create a PowerPoint presentation on the topic "Applications of PowerPoint" with 5 slides. After that, do the following tasks:\pts{10 Marks}
\begin{parts}
  \item Apply a design theme of your choice.
  \item Insert a title and subtitle on the first slide with your name and course.
  \item Apply slide transition effects to all slides.
  \item Insert a shape and animate it to move in a circular path.
  \item Insert an image and add multiple animations to the same image (Entrance + Emphasis + Exit).
  \item Set up your presentation to run automatically with timings for each slide.
  \item Add speaker notes to at least two slides.
\end{parts}

\bigskip
\noindent\textbf{Question 4.} Create the following tables in MS-Access with appropriate data types and constraints:\pts{10 Marks}\\
\textbf{Students}(\underline{StudentID}, Name, Gender, Phone, CourseID)\\
\textbf{Courses}(\underline{CourseID}, CourseName, Duration, Fees)
\begin{parts}
  \item Set StudentID and CourseID as primary keys.
  \item Create a one-to-many relationship between Courses $\rightarrow$ Students.
  \item Enter at least 5 student records and 3 course records.
\end{parts}

\vfill\rule{\linewidth}{0.4pt}
\end{document}
